\documentclass[../main.tex]{subfiles}

\begin{document}

\chapter{可数性公理和分离公理}\label{chapter:ch04-可数性公理和分离公理}

我们现在打算介绍的概念，不像紧致性和连通性那样自然地源于对微积分和分析学的研究，而是出于对拓扑学本身深入研究的需要。例如将一个空间嵌入到度量空间或紧致的 Hausdorff 空间等问题，本质上都是拓扑学的问题，而非分析学的问题。这些特定问题的解决就涉及到可数性公理与分离性公理。

当我们在第 21 节中研究与收敛序列相关的问题时，曾介绍过第一可数性公理。我们也研究过一个分离性公理，即 Hausdorff 公理，并提到过另一个，即 $T_{1}$ 公理。本章我们将介绍另外一些更强的公理，并且研究它们之间的一些关系。主要目的是证明 Urysohn 度量化定理：如果一个拓扑空间 $X$ 满足某种可数性公理（第二可数性公理）和某种分离性公理（正则性公理），那么 $X$ 可以嵌入到某一个度量空间中，因此是可度量化的。

本章最后一节还要介绍几何学上一个十分重要的嵌入定理，即给定一个紧致流形（高维曲面的模拟），我们将证明它可以嵌入到某个有限维的欧氏空间之中。

\subfile{Chapter-04/section-30.tex}
\subfile{Chapter-04/section-31.tex}
\subfile{Chapter-04/section-32.tex}
\subfile{Chapter-04/section-33.tex}
\subfile{Chapter-04/section-34.tex}
\subfile{Chapter-04/section-35.tex}
\subfile{Chapter-04/section-36.tex}
\subfile{Chapter-04/section-extra.tex}

\end{document}